% figuras/tikz/cap07-embudo-vs-bucle.tex
% Embudo lineal (AIDA) frente a recorrido en bucle (consumer decision journey / messy middle).
\begin{tikzpicture}[
    et/.style={rectangle, rounded corners, draw=guindaIPN, thick, fill=guindaIPNClara,
               minimum width=1.9cm, minimum height=0.8cm, align=center, font=\scriptsize},
    flecha/.style={-{Stealth[length=2.2mm]}, thick, draw=doradoIPN},
    retro/.style={-{Stealth[length=2.2mm]}, thick, draw=terracotaIPN, dashed},
    tit/.style={font=\small\bfseries, align=center}
  ]
  % --- Izquierda: embudo lineal
  \node[tit] at (0,3.7) {Embudo lineal};
  \node[et, minimum width=4.6cm] (a1) at (0,2.8) {Atención};
  \node[et, minimum width=3.7cm] (a2) at (0,1.7) {Interés};
  \node[et, minimum width=2.8cm] (a3) at (0,0.6) {Deseo};
  \node[et, minimum width=1.9cm] (a4) at (0,-0.5) {Acción};
  \foreach \a/\b in {a1/a2,a2/a3,a3/a4} { \draw[flecha] (\a) -- (\b); }
  \node[font=\tiny, text=grisIPN, align=center] at (0,-1.4) {una sola dirección;\\ el consumidor solo \enquote{baja}};

  % --- Derecha: bucle
  \node[tit] at (7.2,3.7) {Recorrido en bucle};
  \node[et] (b0) at (7.2,2.8) {Disparador};
  \node[et] (b1) at (5.4,1.1) {Exploración\\ (suma opciones)};
  \node[et] (b2) at (9.0,1.1) {Evaluación\\ (reduce opciones)};
  \node[et] (b3) at (7.2,-0.5) {Compra};
  \node[et] (b4) at (10.6,-0.5) {Uso y\\ reseña};
  \draw[flecha] (b0) -- (b1);
  \draw[retro] (b1.east) to[bend left=25] (b2.west);
  \draw[retro] (b2.west) to[bend left=25] (b1.east);
  \draw[flecha] (b2) -- (b3);
  \draw[flecha] (b3) -- (b4);
  \draw[flecha, draw=bronceIPN, dashed] (b4.north) to[bend right=35] (b1.north);
  \node[font=\tiny, text=grisIPN, align=center] at (7.2,-1.4) {se explora y se evalúa varias veces;\\ la reseña alimenta el recorrido de otros};
\end{tikzpicture}
